\begin{figure}[t]
    \centering
    \includegraphics[width=\linewidth]{images/feature_importance_sql_features_time_rdbms_vs_qiskit.pdf}
    \caption{Importance of various SQL features for SVM and XGBoost trained for optimizing execution time.}
    \label{fig:rdbmsqiskitimportancetime}
\end{figure}
\begin{figure}[t]
    \centering
    \includegraphics[width=\linewidth]{images/feature_importance_sql_features_memory_rdbms_vs_qiskit.pdf}
    \caption{Importance of various SQL features for SVM and RF trained for optimizing memory.}
    \label{fig:rdbmsqiskitimportancemem}
\end{figure}
\section{Additional Results on SQL-only features.}
\label{ssec:sqlf}
Figures~\ref{fig:rdbmsqiskitimportancetime} and~\ref{fig:rdbmsqiskitimportancemem} show that the
dominant SQL signals are structural counts, in particular the number of \texttt{JOIN}s,
predicates, and aggregates. For runtime, the Linear SVM assigns negative coefficients to
\feat{num_where_clauses}, \feat{num_eq_predicates}, and \feat{num_agg_funcs}, suggesting that
queries with heavier filtering and aggregation tend to favor Qiskit Aer. In contrast,
\feat{num_joins} and \feat{num_and_clauses} shift the decision toward the RDBMS backend. For memory,
\feat{num_joins} is the dominant factor, while predicate and aggregation features retain
negative influence, reflecting the benefit of the optimized join algorithms and cost-based planning employed by RDBMS.